\section{Scope Boundaries}
\label{sec:scope}

This section defines explicit boundaries for the Mermaid-superset diagram compiler.

\subsection{Governing Principle}

The scope boundary is defined by the \emph{Mermaid-superset contract}: this compiler accepts
all valid Mermaid diagrams and our extensions, and produces deterministic visual output.
It does not execute, simulate, schedule, or interactively manipulate diagrams.

\begin{quote}
\emph{If Mermaid renders it, we render it (better). Beyond Mermaid, we add composition,
rich theming, animation, and agent IR. We do not add interactivity, data management,
or operational semantics.}
\end{quote}

\subsection{In Scope}

\subsubsection{Diagram Coverage}

All 22 Mermaid diagram types (Section~\ref{sec:mermaid-compat}), organized into five families,
plus our superset extensions (composition, animation, cross-references).

\subsubsection{Input Paths}

\begin{description}
  \item[Mermaid-superset DSL] Text syntax compatible with Mermaid plus our extensions.
  \item[Structured IR] JSON/YAML conforming to per-grammar Domain IR schemas.
\end{description}

\subsubsection{Output Formats}

\begin{description}
  \item[SVG] Primary vector output (source-of-truth backend).
  \item[PNG] Rasterized output via resvg/Skia.
  \item[PDF] Print-quality vector output for documents and decks.
\end{description}

\subsubsection{Key Capabilities}

\begin{description}
  \item[Theming] Complete visual customization via theme files.
  \item[Determinism] Byte-identical output guarantee.
  \item[Composition] Multi-diagram posters and grid layouts.
  \item[Animation] Declarative SMIL animation in SVG output.
  \item[Agent integration] MCP server, CLI, library API.
  \item[Schema validation] Per-grammar JSON Schema validation with clear diagnostics.
\end{description}

\subsection{Out of Scope}

\begin{description}
  \item[Interactive output] No clickable regions, tooltips, or drill-down. Output is static.
  \item[WYSIWYG editing] No drag-and-drop authoring. This is a compiler, not an editor.
  \item[Data management] No persistent storage, collaboration, or version history.
  \item[Operational semantics] No dependency scheduling, resource leveling, critical-path
        analysis, or simulation.
  \item[General data visualization] Not competing with D3/Plotly for arbitrary charts.
        The chart family covers Mermaid's four chart types only.
  \item[Native integrations] Does not natively read from Jira, Azure DevOps, or other
        systems. Consumes only its IR format or Mermaid text.
  \item[Live collaboration] No real-time multi-user editing.
\end{description}

\subsection{Scope Summary}

\begin{table}[h]
\centering
\caption{Scope Boundaries Summary}
\begin{tabularx}{\textwidth}{lX}
\toprule
\textbf{Category} & \textbf{Status} \\
\midrule
All 22 Mermaid diagram types & In Scope \\
Mermaid-superset DSL parsing & In Scope \\
Structured IR (agent path) & In Scope \\
Beautiful themed output (SVG, PNG, PDF) & In Scope \\
Deterministic byte-identical rendering & In Scope \\
Multi-diagram composition/posters & In Scope \\
Declarative animation & In Scope \\
CLI + library + MCP server & In Scope \\
\midrule
Interactive/clickable output & Out of Scope \\
WYSIWYG/canvas editing & Out of Scope \\
Operational/scheduling semantics & Out of Scope \\
Data storage and collaboration & Out of Scope \\
Native data source integrations & Out of Scope \\
General-purpose data visualization & Out of Scope \\
\bottomrule
\end{tabularx}
\end{table}
